\subsection{Fresh-seed comparison at the smaller audit budget}
After the original confirmation exposed the empirical-Bernstein linear floor,
we froze a separate protocol before running seeds 63000--63015. Both the
betting and original EB gates use the same ordinary starts, fitting data,
fresh 4,096-label audit blocks, residual/random candidate proposals, and
branch training. This adds 128 task--seed--gate configurations, 256 controller
trajectories, and 2,048 opportunities. It is a new-seed follow-up motivated by
an earlier finding, not a retroactive change to the original confirmation.
The first experiment runner and all original results remain unchanged.

\begin{table}[tb]
\centering\scriptsize
\resizebox{\linewidth}{!}{\begin{tabular}{llrrrr}
\toprule
Task & Proposal & EB risk & Betting risk & Paired difference [95\% CI] & Width (EB / bet)\\
\midrule
null & residual & 0.25312 & 0.25312 & +0.00000 [+0.00000, +0.00000] & 2.00 / 2.00\\
null & random & 0.25312 & 0.25312 & +0.00000 [+0.00000, +0.00000] & 2.00 / 2.00\\
linear & residual & 0.17210 & 0.17210 & +0.00000 [+0.00000, +0.00000] & 2.00 / 2.00\\
linear & random & 0.17210 & 0.17210 & +0.00000 [+0.00000, +0.00000] & 2.00 / 2.00\\
additive & residual & 0.14698 & 0.13836 & -0.00861 [-0.01078, -0.00644] & 2.00 / 3.44\\
additive & random & 0.14698 & 0.13915 & -0.00782 [-0.01068, -0.00496] & 2.00 / 3.25\\
interaction & residual & 0.17558 & 0.15209 & -0.02349 [-0.02628, -0.02070] & 3.00 / 5.31\\
interaction & random & 0.17357 & 0.14955 & -0.02403 [-0.02629, -0.02177] & 3.00 / 5.56\\
\bottomrule
\end{tabular}
}
\caption{New-seed gate comparison, 16 paired repetitions per task and proposal.
Both gates use 4,096 fresh audit labels per opportunity. Paired difference is
betting minus EB Brier risk, with a descriptive 95\% Student-$t$ interval.
Width is the mean final hidden-unit count. The betting gate's extra products
are not included in training parameter--example matching.}
\label{tab:betting}
\end{table}

The alternative gate provides a substantial constructive improvement for
both proposal mechanisms. For the residual method, additive risk decreases
from 0.14698 under EB to 0.13836 under betting, a paired change of
$-0.00861$ with interval $[-0.01078,-0.00644]$. Interaction risk decreases
from 0.17558 to 0.15209, a paired change of $-0.02349$ with interval
$[-0.02628,-0.02070]$. Mean widths increase from 2 to 3.4375 on additive
data and from 3 to 5.3125 on interaction data. Null and linear networks
remain at width two with identical predictions under both gates.
These are repeated trained histories with a finite-sample commitment
guarantee, rather than forced-growth examples.

The improvement is at the same audit-label budget, not the same final
training effort. Accepted wider models require more later branch training:
residual training/probe parameter--example work rises by factors 1.206 and
1.224 on additive and interaction data. The betting calculations add
$15\cdot4096\cdot8=491{,}520$ factor evaluations per trajectory, including
rejected comparisons. Current research instrumentation also computes EB
statistics in both arms. Both consume 32,768 audit labels plus 2,048 fitting
labels; no evidence-sharing gain from Theorem~\ref{thm:shared-evidence} is
claimed for this restarted testing procedure.

The follow-up also sharpens the adverse selector evidence. Under the betting
gate, residual minus random is $-0.000789$ with interval
$[-0.003284,0.001705]$ on additive data and $+0.002542$ with interval
$[0.000242,0.004841]$ on interaction data. Random thus retains a favorable
interaction result while avoiding candidate optimization. Comparisons of
the residual method with its retrospectively matched fixed capacity and
training/probe budget remain unresolved on both nonlinear tasks. The random
betting method has a modest favorable additive comparison against its own
matched fixed model, $-0.00194$ with interval $[-0.00353,-0.00035]$; that
positive control result is retained. Larger full-label-pool networks still
have lower mean risks, with the capacity and offline-optimization caveats
described earlier.

Independent checks reconstruct all 1,024 betting-arm opportunities and
3,072 product tests from saved parameters and regenerated audit examples;
maximum log-e discrepancy is $4.0\cdot10^{-13}$. Finite-binomial null
calculations check the expectation and tail bound directly. No betting-arm
retained step increases evaluator risk in this follow-up, but that finite
observation is not the proof of validity. The proof is
Proposition~\ref{prop:betting-gate}, and evaluation remains Monte Carlo in
the input distribution.

\begin{figure}[tb]
\centering\includegraphics[width=\linewidth]{betting_followup.pdf}
\caption{Measured, fresh-seed gate follow-up at the same smaller audit budget.
Paired intervals and trajectories show that a less conservative established
test can enable repeated growth. Extra accepted capacity raises subsequent
training work. Both informative and random proposals benefit, so the result
attributes progress to the evidence mechanism rather than validating the
residual proposal score.}
\label{fig:betting}
\end{figure}
